% ============================================================
% f2014_shandong_science_q17_prism.tex
% 2014 山东卷（理科）第 17 题：四棱柱 ABCD-A1B1C1D1
% 来源 PDF：普通高考/2014/2014山东理.pdf
% 原页：第 1 页（A3 横排）右栏，印面 907，第 17 题题干下方
% 结构：底面 ABCD 为等腰梯形，AB=2CD=2，∠DAB=60°，M 为 AB 中点。
%       实线：AB、AA1、A1D1、D1C1、C1B1、BB1、A1B1；
%       虚线（不可见棱）：AD、DC、CB、DD1、CC1；
%       虚线（辅助线）：D1A、D1M、C1M。
% 比例：AB:CD=2:1 按屏幕长 4.6:2.3 画出，侧棱方向取 (0.95,2.85)。
% ============================================================
\begin{tikzpicture}[>=Stealth, line width=0.7pt, font=\small,
  x=1cm, y=1cm, line cap=round, line join=round]
  % 底面：AB=4.6（屏幕长），CD=2.3，梯形
  \coordinate (A)  at (0.10,0.30);
  \coordinate (B)  at (4.70,0.30);
  \coordinate (M)  at (2.40,0.30);
  \coordinate (D)  at (1.45,2.05);
  \coordinate (C)  at (3.75,2.05);
  % 侧棱平移向量 (0.95,2.85)
  \coordinate (A1) at (1.05,3.15);
  \coordinate (B1) at (5.65,3.15);
  \coordinate (D1) at (2.40,4.90);
  \coordinate (C1) at (4.70,4.90);
  % 不可见棱
  \draw[dashed] (A) -- (D);
  \draw[dashed] (D) -- (C);
  \draw[dashed] (C) -- (B);
  \draw[dashed] (D) -- (D1);
  \draw[dashed] (C) -- (C1);
  % 辅助线
  \draw[dashed] (D1) -- (A);
  \draw[dashed] (D1) -- (M);
  \draw[dashed] (C1) -- (M);
  % 可见棱
  \draw (A) -- (B);
  \draw (A) -- (A1);
  \draw (A1) -- (D1);
  \draw (D1) -- (C1);
  \draw (C1) -- (B1);
  \draw (B) -- (B1);
  \draw (A1) -- (B1);
  % 顶点与中点标注
  \node[below left]  at (A)  {$A$};
  \node[below right] at (B)  {$B$};
  \node[below]       at (M)  {$M$};
  \node at (1.10,1.94) {$D$};
  \node at (3.98,1.84) {$C$};
  \node[left]        at (A1) {$A_1$};
  \node[right]       at (B1) {$B_1$};
  \node[above left]  at (D1) {$D_1$};
  \node[above right] at (C1) {$C_1$};
\end{tikzpicture}
